% Options for packages loaded elsewhere
\PassOptionsToPackage{unicode}{hyperref}
\PassOptionsToPackage{hyphens}{url}
\PassOptionsToPackage{space}{xeCJK}
\documentclass[
]{article}
\usepackage{xcolor}
\usepackage{amsmath,amssymb}
\setcounter{secnumdepth}{-\maxdimen} % remove section numbering
\usepackage{iftex}
\ifPDFTeX
  \usepackage[T1]{fontenc}
  \usepackage[utf8]{inputenc}
  \usepackage{textcomp} % provide euro and other symbols
\else % if luatex or xetex
  \usepackage{unicode-math} % this also loads fontspec
  \defaultfontfeatures{Scale=MatchLowercase}
  \defaultfontfeatures[\rmfamily]{Ligatures=TeX,Scale=1}
\fi
\usepackage{lmodern}
\ifPDFTeX\else
  % xetex/luatex font selection
  \setmainfont[]{Times New Roman}
  \setsansfont[]{SimHei}
  \setmonofont[]{SimKai}
  \ifXeTeX
    \usepackage{xeCJK}
    \setCJKmainfont[]{SimSun}
  \fi
  \ifLuaTeX
    \usepackage[]{luatexja-fontspec}
    \setmainjfont[]{SimSun}
  \fi
\fi
% Use upquote if available, for straight quotes in verbatim environments
\IfFileExists{upquote.sty}{\usepackage{upquote}}{}
\IfFileExists{microtype.sty}{% use microtype if available
  \usepackage[]{microtype}
  \UseMicrotypeSet[protrusion]{basicmath} % disable protrusion for tt fonts
}{}
\makeatletter
\@ifundefined{KOMAClassName}{% if non-KOMA class
  \IfFileExists{parskip.sty}{%
    \usepackage{parskip}
  }{% else
    \setlength{\parindent}{0pt}
    \setlength{\parskip}{6pt plus 2pt minus 1pt}}
}{% if KOMA class
  \KOMAoptions{parskip=half}}
\makeatother
\usepackage{longtable,booktabs,array}
\newcounter{none} % for unnumbered tables
\usepackage{calc} % for calculating minipage widths
% Correct order of tables after \paragraph or \subparagraph
\usepackage{etoolbox}
\makeatletter
\patchcmd\longtable{\par}{\if@noskipsec\mbox{}\fi\par}{}{}
\makeatother
% Allow footnotes in longtable head/foot
\IfFileExists{footnotehyper.sty}{\usepackage{footnotehyper}}{\usepackage{footnote}}
\makesavenoteenv{longtable}
\usepackage{graphicx}
\makeatletter
\newsavebox\pandoc@box
\newcommand*\pandocbounded[1]{% scales image to fit in text height/width
  \sbox\pandoc@box{#1}%
  \Gscale@div\@tempa{\textheight}{\dimexpr\ht\pandoc@box+\dp\pandoc@box\relax}%
  \Gscale@div\@tempb{\linewidth}{\wd\pandoc@box}%
  \ifdim\@tempb\p@<\@tempa\p@\let\@tempa\@tempb\fi% select the smaller of both
  \ifdim\@tempa\p@<\p@\scalebox{\@tempa}{\usebox\pandoc@box}%
  \else\usebox{\pandoc@box}%
  \fi%
}
% Set default figure placement to htbp
\def\fps@figure{htbp}
\makeatother
\setlength{\emergencystretch}{3em} % prevent overfull lines
\providecommand{\tightlist}{%
  \setlength{\itemsep}{0pt}\setlength{\parskip}{0pt}}
\usepackage{bookmark}
\IfFileExists{xurl.sty}{\usepackage{xurl}}{} % add URL line breaks if available
\urlstyle{same}
\hypersetup{
  hidelinks,
  pdfcreator={LaTeX via pandoc}}

\author{}
\date{}

\begin{document}

分类号 密级

U D C 编号

\pandocbounded{\includegraphics[keepaspectratio]{media/image1.jpeg}}
\pandocbounded{\includegraphics[keepaspectratio]{media/image2.png}}

本科毕业论文(设计)

\begin{quote}
\textbf{题目}

学 院

专 业 名 称

年 级

学 生 姓 名

学 号

指 导 教 师 王海军；李云雪
\end{quote}

二〇二五 年 四 月

\subsubsection{}\label{section}

\textbf{湖北第二师范学院毕业论文（设计）原创性声明}

本人郑重声明：所呈交的学位论文，是本人在导师的指导下，独立进行研究工作所取得的成果。除文中已经注明引用的内容外，本论文不含任何其他个人或集体已经发表或撰写过的作品成果。对本文的研究做出重要贡献的个人和集体，均已在文中以明确方式标明。本人完全意识到本声明的法律结果由本人承担。

论文作者签名： 年 月 日

\textbf{摘要：}文本文本文本文本文本文本文本文本文本文本文本文本文本文本文本文本文本文本文本文本文本文本文本文本文本文本文本文本文本文本文本文本文本文本文本文本文本文本文本文本文本文本文本文本文本文本文本文本文本文本文本文本文本文本文本。

文本文本文本文本文本文本文本文本文本文本文本文本文本文本文本文本文本文本文本文本文本文本文本文本文本文本文本文本文本文本文本文本文本文本文本文本文本文本文本文本。文本文本文本文本文本文本文本文本文本文本文本文本文本文本文本文本文本文本文本文本文本文本文本文本文本文本文本文本文本文本文本文本文本文本文本文本文本文本文本文本。

\textbf{关键词：}关键词1；关键词2；关键词3；关键词4

\textbf{Abstract:} text text text text text text text text text text
text text text text text, text text text text text text text, text text
text text text text text text text text text text text text. Text text
text text text text text text text text text text text text text, text
text text text text text text, text text text text text text text text
text text text text text text. Text text text text text text text text
text text text text text text text, text text text text text text text,
text text text text text text text text text text text text text text.
Text text text text text text text text text text text text text text
text, text text text text text text text, text text text text text text
text text text text text text text text. Text text text text text text
text text text text text text text text text, text text text text text
text text, text text text text text text text text text text text text
text text.

Second paragraph text text text text text text text text text text text
text text text, text text text text text text text, text text text text
text text text text text text text text text text. Text text text text
text text text text text text text text text text text, text text text
text text. Second paragraph text text text text text text text text text
text text text text text, text text text text text text text, text text
text text text text text text text text text text text text. Text text
text text text text text text text text text text text text text, text
text text text text.

\textbf{Keywords:} keywords1; keywords2; keywords3; keywords4

\textbf{目录}

\hyperref[_Toc15630]{1. 绪论 \hyperref[_Toc15630]{1}}

\hyperref[_Toc28684]{1.1 标题 \hyperref[_Toc28684]{1}}

\hyperref[_Toc5342]{1.1.1 标题 \hyperref[_Toc5342]{2}}

\hyperref[_Toc26435]{1.1.2 标题 \hyperref[_Toc26435]{4}}

\hyperref[_Toc32459]{1.1.3 标题 \hyperref[_Toc32459]{4}}

\hyperref[_Toc22971]{1.2 标题 \hyperref[_Toc22971]{5}}

\hyperref[_Toc7083]{1.3 标题 \hyperref[_Toc7083]{5}}

\hyperref[_Toc31455]{2. 标题 \hyperref[_Toc31455]{6}}

\hyperref[_Toc4287]{2.1 标题 \hyperref[_Toc4287]{6}}

\hyperref[_Toc12161_WPSOffice_Level3]{2.1.1 标题
\hyperref[_Toc12161_WPSOffice_Level3]{6}}

\hyperref[_Toc18030]{3. 标题 \hyperref[_Toc18030]{7}}

\hyperref[_Toc7093_WPSOffice_Level2]{3.1 标题
\hyperref[_Toc7093_WPSOffice_Level2]{7}}

\hyperref[_Toc7714]{3.1.1 标题 \hyperref[_Toc7714]{7}}

\hyperref[_Toc4937]{4. 标题 \hyperref[_Toc4937]{8}}

\hyperref[_Toc22371]{4.1 标题 \hyperref[_Toc22371]{8}}

\hyperref[_Toc27652]{4.1.1 标题 \hyperref[_Toc27652]{8}}

\hyperref[_Toc4867]{5. 结论 \hyperref[_Toc4867]{9}}

\hyperref[_Toc26992]{参考文献 \hyperref[_Toc26992]{10}}

\hyperref[_Toc9629]{致谢 \hyperref[_Toc9629]{12}}

\hyperref[_Toc32605]{附录 \hyperref[_Toc32605]{13}}

\protect\phantomsection\label{_Toc15630}{}\textbf{1. 绪论}

正文文字正文文字正文文字正文文字正文文字正文文字正文文字正文文字正文文字正文文字正文文字正文文字正文文字正文文字正文文字正文文字正文文字正文文字正文文字正文文字正文文字正文文字正文文字正文文字\textsuperscript{{[}1{]}}。

\protect\phantomsection\label{_Toc28684}{}\textbf{1.1 标题}

正文文字正文文字正文文字正文文字正文文字正文文字正文文字正文文字正文文正文文字正文文字正文文字正文文字正文文字正文文字\textsuperscript{{[}2{]}}。

正文文字正文文字正文文字正文文字正文文字正文文字正文文字正文文字正文文正文文字正文文字正文文字正文文字正文文字正文文字。

正文文字正文文字正文文字正文文字正文文字正文文字正文文字正文文字正文文正文文字正文文字正文文字正文文字正文文字正文文字\textsuperscript{{[}3{]}}。

正文文字正文文字正文文字正文文字正文文字正文文字正文文字正文文字正文文正文文字正文文字正文文字正文文字正文文字正文文字。

表1-1 表名

{\def\LTcaptype{none} % do not increment counter
\begin{longtable}[]{@{}
  >{\raggedright\arraybackslash}p{(\linewidth - 8\tabcolsep) * \real{0.2622}}
  >{\raggedright\arraybackslash}p{(\linewidth - 8\tabcolsep) * \real{0.2965}}
  >{\raggedright\arraybackslash}p{(\linewidth - 8\tabcolsep) * \real{0.2130}}
  >{\raggedright\arraybackslash}p{(\linewidth - 8\tabcolsep) * \real{0.0736}}
  >{\raggedright\arraybackslash}p{(\linewidth - 8\tabcolsep) * \real{0.0809}}@{}}
\toprule\noalign{}
\endhead
\bottomrule\noalign{}
\endlastfoot
名称 & 字段名称 & 数据类型 & 主键 & 非空 \\
公司编号 & COM\_ID & VARCHAR(30) & YES & YES \\
月份 & BAL\_MONTH & CHAR (6) & YES & YES \\
商品编号 & ITEM\_ID & VARCHAR(30) & YES & YES \\
当月销售数量 & QTY\_SOLD & NUMERIC(18,2) & NO & NO \\
当月销售含税金额 & AMT\_SOLD\_WITH\_TAX & NUMERIC (18,2) & NO & NO \\
当月毛利 & GROSS\_PROFIT & NUMERIC (18,2) & NO & NO \\
\end{longtable}
}

注：如有需要可对表格进行注释说明，可省略。

数据来源：如需对数据来源进行说明，请参照此格式，可省略。

正文文字正文文字正文文字正文文字正文文字正文文字正文文字正文文字正文文字正文文字正文文字正文文字正文文字正文文字正文文字正文文字正文文字正文文字正文文字正文文字正文文字正文文字正文文字正文文字正文文字正文文字。

\protect\phantomsection\label{_Toc5342}{}\textbf{1.1.1 标题}

正文文字正文文字正文文字正文文字正文文字正文文字正文文字正文文字正文文字正文文字正文文字正文文字正文文字正文文字正文文字正文文字正文文字正文文字正文文字正文文字正文文字正文文字正文文字正文文字。

s\emph{in}α-s\emph{in}β=2s\emph{in}\pandocbounded{\includegraphics[keepaspectratio]{media/image3.wmf}}\emph{co}s\pandocbounded{\includegraphics[keepaspectratio]{media/image4.wmf}}
（1.1）

正文文字正文文字正文文字正文文字正文文字正文文字正文文字正文文字正文文字正文文字正文文字正文文字正文文字正文文字正文文字正文文字正文文字正文文字正文文字正文文字正文文字正文文字正文文字正文文字。

\pandocbounded{\includegraphics[keepaspectratio]{media/image5.png}}

图1-1 图名

注：如有需要可对图片进行注释说明，可省略。

图片来源：如需对图片来源进行说明，请参照此格式，可省略。

正文文字正文文字正文文字正文文字正文文字正文文字正文文字正文文字正文文字正文文字正文文字正文文字正文文字正文文字正文文字正文文字正文文字正文文字正文文字正文文字正文文字正文文字正文文字正文文字。

正文文字正文文字正文文字正文文字正文文字正文文字正文文字正文文字正文文字正文文字正文文字正文文字正文文字正文文字正文文字正文文字正文文字正文文字正文文字正文文字正文文字正文文字正文文字正文文字。

正文文字正文文字正文文字正文文字正文文字正文文字正文文字正文文字正文文字正文文字正文文字正文文字正文文字正文文字正文文字正文文字正文文字正文文字正文文字正文文字正文文字正文文字正文文字正文文字正文文字正文文字正文文字正文文字正文文字正文文字正文文字正文文字正文文字正文文字正文文字正文文字正文文字。

\pandocbounded{\includegraphics[keepaspectratio]{media/image6.png}}

(a) 分图1

\pandocbounded{\includegraphics[keepaspectratio]{media/image7.png}}

(b) 分图2

图1-2 图名

正文文字正文文字正文文字正文文字正文文字正文文字正文文字正文文字正文文字正文文字正文文字正文文字正文文字正文文字正文文字正文文字正文文字正文文字正文文字正文文字正文文字正文文字正文文字正文文字。

正文文字正文文字正文文字正文文字正文文字正文文字正文文字正文文字正文文字正文文字正文文字正文文字正文文字正文文字正文文字正文文字正文文字正文文字正文文字正文文字正文文字正文文字正文文字正文文字。文文字正文文字正文文字正文文字正文文字正文文字正文文字正文文字正文文字正文文字正文文字正文文字正文文字正文文字正文文字正文文字正文文字正文文字正文文字正文文字正文文字正文文字正文文字正文文字。

正文文字正文文字正文文字正文文字正文文字正文文字正文文字正文文字正文文字正文文字正文文字正文文字正文文字正文文字正文文字正文文字正文文字正文文字正文文字正文文字正文文字正文文字正文文字正文文字正文文字正文文字正文文字。

\protect\phantomsection\label{_Toc26435}{}\textbf{1.1.2 标题}

正文文字正文文字正文文字正文文字正文文字正文文字正文文字正文文字正文文字正文文字正文文字正文文字正文文字正文文字正文文字正文文字正文文字正文文字正文文字正文文字正文文字正文文字正文文字正文文字。

正文文字正文文字正文文字正文文字正文文字正文文字正文文字正文文字正文文字正文文字正文文字正文文字正文文字正文文字正文文字正文文字正文文字正文文字正文文字正文文字正文文字正文文字正文文字正文文字。

\protect\phantomsection\label{_Toc32459}{}\textbf{1.1.3 标题}

正文文字正文文字正文文字正文文字正文文字正文文字正文文字正文文字正文文字正文文字正文文字正文文字正文文字正文文字正文文字正文文字正文文字正文文字正文文字正文文字正文文字正文文字正文文字正文文字正文文字正文文字正文文字正文文字正文文字正文文字正文文字正文文字正文文字正文文字正文文字正文文字正文文字正文文字正文文字正文文字正文文字正文文字正文文字正文文字正文文字正文文字正文文字正文文字正文文字。

表1-2 表名

{\def\LTcaptype{none} % do not increment counter
\begin{longtable}[]{@{}
  >{\raggedright\arraybackslash}p{(\linewidth - 8\tabcolsep) * \real{0.1770}}
  >{\raggedright\arraybackslash}p{(\linewidth - 8\tabcolsep) * \real{0.1770}}
  >{\raggedright\arraybackslash}p{(\linewidth - 8\tabcolsep) * \real{0.1770}}
  >{\raggedright\arraybackslash}p{(\linewidth - 8\tabcolsep) * \real{0.1770}}
  >{\raggedright\arraybackslash}p{(\linewidth - 8\tabcolsep) * \real{0.1770}}@{}}
\toprule\noalign{}
\endhead
\bottomrule\noalign{}
\endlastfoot
编号 & 长 & 宽 & 高 & 检测结果 \\
10001 & 104 & 98 & 103 & 合格 \\
10002 & 103 & 105 & 97 & 合格 \\
10003 & 101 & 97 & 94 & 不合格 \\
10004 & 106 & 101 & 106 & 不合格 \\
10005 & 98 & 99 & 104 & 合格 \\
10006 & 94 & 102 & 94 & 不合格 \\
10007 & 106 & 95 & 106 & 不合格 \\
10008 & 102 & 99 & 106 & 不合格 \\
\end{longtable}
}

续表1-2 表名

{\def\LTcaptype{none} % do not increment counter
\begin{longtable}[]{@{}
  >{\raggedright\arraybackslash}p{(\linewidth - 8\tabcolsep) * \real{0.1770}}
  >{\raggedright\arraybackslash}p{(\linewidth - 8\tabcolsep) * \real{0.1770}}
  >{\raggedright\arraybackslash}p{(\linewidth - 8\tabcolsep) * \real{0.1770}}
  >{\raggedright\arraybackslash}p{(\linewidth - 8\tabcolsep) * \real{0.1770}}
  >{\raggedright\arraybackslash}p{(\linewidth - 8\tabcolsep) * \real{0.1770}}@{}}
\toprule\noalign{}
\endhead
\bottomrule\noalign{}
\endlastfoot
编号 & 长 & 宽 & 高 & 检测结果 \\
10009 & 99 & 98 & 98 & 合格 \\
10010 & 101 & 97 & 94 & 不合格 \\
10011 & 106 & 101 & 106 & 不合格 \\
10012 & 98 & 99 & 104 & 合格 \\
10013 & 94 & 102 & 94 & 不合格 \\
10014 & 106 & 95 & 106 & 不合格 \\
10015 & 102 & 99 & 106 & 不合格 \\
10016 & 99 & 98 & 98 & 合格 \\
\end{longtable}
}

\protect\phantomsection\label{_Toc22971}{}\textbf{1.2 标题}

正文文字正文文字正文文字正文文字正文文字正文文字正文文字正文文字正文文字正文文字正文文字正文文字正文文字正文文字正文文字正文文字正文文字正文文字正文文字正文文字正文文字正文文字正文文字正文文字。

void COM(void) // 温湿写入

\{

uchar i;

for(i=0;i\textless8;i++)

\{

ucharFLAG=2;

while((!DATA\_PIN)\&\&ucharFLAG++);

Delay\_10us();

uchartemp=0;

if(DATA\_PIN)uchartemp=1;

ucharFLAG=2;

while((DATA\_PIN)\&\&ucharFLAG++);

if(ucharFLAG==1)break;

ucharcomdata\textless\textless=1;

ucharcomdata\textbar=uchartemp;

\}

\}

\protect\phantomsection\label{_Toc7083}{}初始帧图像输入：Q=\{Q\textbar t=1,2,3,\ldots,L\ldots,N\}

确定背景模型

初始化颜色背景：M（p）=\{D\textsubscript{1},
D\textsubscript{2},\ldots D\textsubscript{N}\}

For p=1-\/-\/-\textgreater 高×宽

For k=1-\/-\/-\textgreater m

计算不同时刻颜色相似度 C（Q\textsubscript{1}，Q\textsubscript{k}）

END

END

\textbf{1.3 标题}

正文文字正文文字正文文字正文文字正文文字正文文字正文文字正文文字正文文字正文文字正文文字正文文字正文文字正文文字正文文字正文文字正文文字正文文字正文文字正文文字正文文字正文文字正文文字正文文字正文文字正文文字正文文字正文文字正文文字正文文字正文文字正文文字正文文字正文文字正文文字正文文字正文文字正文文字正文文字。

\begin{enumerate}
\def\labelenumi{\arabic{enumi}.}
\setcounter{enumi}{1}
\item
  \textbf{\hfill\break
  }\protect\phantomsection\label{_Toc31455}{}\textbf{标题}
\end{enumerate}

正文文字正文文字正文文字正文文字正文文字正文文字正文文字正文文字正文文字正文文字正文文字正文文字正文文字正文文字正文文字正文文字正文文字正文文字正文文字正文文字正文文字正文文字正文文字正文文字。

\protect\phantomsection\label{_Toc4287}{}\textbf{2.1 标题}

正文文字正文文字正文文字正文文字正文文字正文文字正文文字正文文字正文文字正文文字正文文字正文文字正文文字正文文字正文文字正文文字正文文字正文文字正文文字正文文字正文文字正文文字正文文字正文文字。

\protect\phantomsection\label{_Toc12161_WPSOffice_Level3}{}\textbf{2.1.1
标题}

正文文字正文文字正文文字正文文字正文文字正文文字正文文字正文文字正文文字正文文字正文文字正文文字正文文字正文文字正文文字正文文字正文文字正文文字正文文字正文文字正文文字正文文字正文文字正文文字。

正文文字正文文字正文文字正文文字正文文字正文文字正文文字正文文字正文文字正文文字正文文字正文文字正文文字正文文字正文文字正文文字正文文字正文文字正文文字正文文字正文文字正文文字正文文字正文文字。

正文文字正文文字正文文字正文文字正文文字正文文字正文文字正文文字正文文字正文文字正文文字正文文字正文文字正文文字正文文字正文文字正文文字正文文字正文文字正文文字正文文字正文文字正文文字正文文字。

\begin{enumerate}
\def\labelenumi{\arabic{enumi}.}
\setcounter{enumi}{2}
\item
  \textbf{\hfill\break
  }\protect\phantomsection\label{_Toc18030}{}\textbf{标题}
\end{enumerate}

正文文字正文文字正文文字正文文字正文文字正文文字正文文字正文文字正文文字正文文字正文文字正文文字正文文字正文文字正文文字正文文字正文文字正文文字正文文字正文文字正文文字正文文字正文文字正文文字。

\protect\phantomsection\label{_Toc7093_WPSOffice_Level2}{}\textbf{3.1
标题}

正文文字正文文字正文文字正文文字正文文字正文文字正文文字正文文字正文文字正文文字正文文字正文文字正文文字正文文字正文文字正文文字正文文字正文文字正文文字正文文字正文文字正文文字正文文字正文文字。

\protect\phantomsection\label{_Toc7714}{}\textbf{3.1.1 标题}

正文文字正文文字正文文字正文文字正文文字正文文字正文文字正文文字正文文字正文文字正文文字正文文字正文文字正文文字正文文字正文文字正文文字正文文字正文文字正文文字正文文字正文文字正文文字正文文字。

正文文字正文文字正文文字正文文字正文文字正文文字正文文字正文文字正文文字正文文字正文文字正文文字正文文字正文文字正文文字正文文字正文文字正文文字正文文字正文文字正文文字正文文字正文文字正文文字正文文字正文文字正文文字正文文字。

正文文字正文文字正文文字正文文字正文文字正文文字正文文字正文文字正文文字正文文字正文文字正文文字正文文字正文文字正文文字正文文字正文文字正文文字正文文字正文文字正文文字正文文字正文文字正文文字正文文字正文文字正文文字正文文字正文文字正文文字正文文字正文文字正文文字正文文字正文文字正文文字。

\textbf{\hfill\break
}\protect\phantomsection\label{_Toc4937}{}\textbf{4. 标题}

正文文字正文文字正文文字正文文字正文文字正文文字正文文字正文文字正文文字正文文字正文文字正文文字正文文字正文文字正文文字正文文字正文文字正文文字正文文字正文文字正文文字正文文字正文文字正文文字。

\protect\phantomsection\label{_Toc22371}{}\textbf{4.1 标题}

正文文字正文文字正文文字正文文字正文文字正文文字正文文字正文文字正文文字正文文字正文文字正文文字正文文字正文文字正文文字正文文字正文文字正文文字正文文字正文文字正文文字正文文字正文文字正文文字。

\protect\phantomsection\label{_Toc27652}{}\textbf{4.1.1 标题}

正文文字正文文字正文文字正文文字正文文字正文文字正文文字正文文字正文文字正文文字正文文字正文文字正文文字正文文字正文文字正文文字正文文字正文文字正文文字正文文字正文文字正文文字正文文字正文文字。

正文文字正文文字正文文字正文文字正文文字正文文字正文文字正文文字正文文字正文文字正文文字正文文字正文文字正文文字正文文字正文文字正文文字正文文字正文文字正文文字正文文字正文文字正文文字正文文字。

正文文字正文文字正文文字正文文字正文文字正文文字正文文字正文文字正文文字正文文字正文文字正文文字正文文字正文文字正文文字正文文字正文文字正文文字正文文字正文文字正文文字正文文字正文文字正文文字。

\textbf{\hfill\break
}\protect\phantomsection\label{_Toc4867}{}\textbf{5. 结论}

正文文字正文文字正文文字正文文字正文文字正文文字正文文字正文文字正文文字正文文字正文文字正文文字正文文字正文文字正文文字正文文字正文文字正文文字正文文字正文文字正文文字正文文字正文文字正文文字。

正文文字正文文字正文文字正文文字正文文字正文文字正文文字正文文字正文文字正文文字正文文字正文文字正文文字正文文字正文文字正文文字正文文字正文文字正文文字正文文字正文文字正文文字正文文字正文文字。

\textbf{\hfill\break
}\protect\phantomsection\label{_Toc26992}{}\textbf{参考文献}

{[}1{]} 杨亚男, 付春玲. 视频目标跟踪算法综述{[}J{]}. 科技资讯, 2018,
16(02): 14-15.

{[}2{]} 徐群, 李家辉, 陈琛, 等.
人工智能技术的智能视频监控系统的相关探究{[}J{]}. 山东工业技术, 2019,
12(09): 131-134.

{[}3{]} 杨艳. 智能交通中运动目标检测与跟踪算法研究{[}D{]}. 昆明:
昆明理工大学, 2016.

{[}4{]} 郜璐璐. 浅析基于多目标跟踪的医学图像分析{[}J{]}. 影像技术. 2014,
29(3): 35-36.

{[}5{]} 张剑华. 均值漂移算法在视频运动目标跟踪中的应用研究{[}D{]}. 西安:
西安电子科技大学, 2015.

{[}6{]} 周梦姣. 基于多特征融合的目标跟踪算法{[}D{]}. 西安:
西安电子科技大学, 2017.

{[}7{]} Han Z, Zhang R, L Wen, et al. Moving Object Tracking Method
Based on Improved Camshift Algorithm{[}J{]}. International Conference on
Industrial Informatics-computing Technology, 2016, pp: 91-95.

{[}8{]} Guo XJ, Wan L, Liu F, et al. An Algorithm for Small Space Target
Tracking Based on Extended Kalman Filter{[}J{]}. Electronics Optics \&
Control, 2016, 4(2): 36-41.

{[}9{]} 权义萍, 金鑫, 张蕾, 等. 基于 Mean-Shift
的卡尔曼粒子滤波车辆跟踪算法{[}J{]}. 南京: 南京工业大学学报, 2014, 3(1):
12-16.

{[}10{]} 宁纪锋. 图像分割和目标跟踪中的若干问题研究{[}D{]}. 西安:
西安电子科技大学, 2009.

{[}11{]} 朱胜利. Mean-Shift 及相关算法在视频跟踪中的研究{[}D{]}. 杭州:
浙江大学, 2006.

{[}12{]} 龙煜. 运动目标检测和遮挡条件下目标跟踪的研究{[}D{]}. 武汉:
武汉科技大学, 2018.

{[}13{]} SHA Musavi, BS Chowdhry, J Bhatti. Object Tracking based on
ActiveContour Modeling{[}C{]}. In: Proceedings of International
Conference on Wireless Communications. 2014, pp: 1-5.

{[}14{]} Weiwu Ren, Xiao Chen, Xiaoming Wang, et al. A Novel Optical
Flow Algorithm Based on Bionic Features for Robust Tracing{[}M{]}.
Germany: Springer International Publishing. 2015, pp: 235-240.

{[}15{]} Dušan Krokavec and Anna Filasová. Fault Residuals Based on
Distributed Discrete-Time Linear Kalman Filtering{[}M{]}. London:
IntechOpen, 2018. pp: 110-125.

{[}16{]} 贺超. 基于图像频域分析显著目标检测算法研究{[}D{]}.
济南:山东大学, 2017.

{[}17{]} 黄金海. 背景帧间差分法的移动目标跟踪研究{[}J{]}. 中国仪器仪表,
2019(01): 62-65.

{[}18{]} 王奎奎, 玉振明. 融合背景减法和帧差法的运动目标检测{[}J{]}.
电视技术, 2015, 39(24): 94-99.

{[}19{]} 胡以静, 李政访, 胡跃明. 基于光流的运动分析理论及应用{[}J{]}.
计算机测量与控制, 2007(02): 219-221.

{[}20{]} 张艳艳. 基于光流算法的运动目标检测应用研究{[}D{]}. 西安:
西安石油大学, 2018.

{[}21{]} Higuchi T. Monte-Carlo filter using the geneic algorithm
operators{[}J{]}. Journal of Statistical Computation and Simulation,
1997, 59(1): 1-23.

{[}22{]} 魏小华, 巫少龙, 叶志斌.
基于改进粒子滤波的视觉目标跟踪算法研究{[}J{]}. 电视术, 2018, 42(11):
72-77.

{[}23{]} Q. Zhou, J. K. Aggarwal. Object Tracking in an Outdoor
Environment Using Fusion of Features and Cameras{[}J{]}. Image and
Vision Computing, 2006, 24(18): 1244--1255.

{[}24{]} 程泽丰, 潘翔. 序贯贝叶斯滤波估计模函数和水平波数{[}J{]}.
杭州电子科技大学学报(自然科学版), 2017, 37(01): 46-50.

{[}25{]} 白薇. 拟蒙特卡洛方法下Heston模型的离散化{[}D{]}. 北京:
清华大学, 2015.

{[}26{]} 赵勇胜, 赵拥军, 赵闯.
利用重要性采样的时差-频差联合估计算法{[}J{]}. 航空报, 2017, 38(01):
195-202.

{[}27{]} 陈潇. 基于特征的视频跟踪算法研究{[}D{]}. 西安:
西安电子科技大学, 2008.

{[}28{]} 黄建. 目标跟踪的粒子滤波算法研究{[}D{]}. 成都: 电子科技大学,
2012.

{[}29{]} Kun Yu, Ting Li, Jie Chen, et al. Classification Method for
Object Feature Extraction Based on Laser Scanning Data{[}M{]}. Germany:
Springer Berlin Heidelberg, 2013, pp: 125-132.

{[}30{]} 陈倩, 陶彬娇, 潘中良, 等.
主颜色提取算法及其在服装图像检索中的应用{[}J{]}.
华南师范大学学报(自然科学版), 2014, 12(1): 1-8.

{[}31{]} 宋玉婷. 基于三维彩色直方图均衡化的彩色图像增强算法研究{[}D{]}.
济南: 山东财经大学, 2013.

{[}32{]} 郝剑南. 纹理增强提取与面向对象结合的高分影像分类的应用{[}D{]}.
南昌: 东华理工大学, 2018.

{[}33{]} 程传阳. 高光谱遥感图像融合与边缘提取方法研究{[}D{]}. 南昌:
南昌航空大学, 2018.

{[}34{]} Hui Fang, Aimin Zhou, Guixu Zhang. A Mean Shift Assisted
Differential Evolution Algorithm{[}M{]}. Germany: Springer: Singapore.
2016.

{[}35{]} 彭志勇, 别秀德. 一种增强型Mean Shift跟踪算法{[}J{]}.
可编程控制器与工厂自动化, 2015(09): 99-102.

{[}36{]} 程佳佳. 基于改进Mean Shift算法的图像分割和目标跟踪{[}D{]}.
天津: 天津大学, 2017.

{[}37{]} Kun Chen, ChunLei Liu, Yongjin Xu. Face Detection and Tracking
Based on Adaboost CamShift and Kalman Filter Algorithm{[}M{]}. Germany:
Springer Berlin Heidelberg. 2014, pp: 45-51.

{[}38{]} 王伟程. Camshift目标跟踪算法改进研究{[}J{]}. 石油化工自动化,
2019, 55(01): 35-38.

{[}39{]} Rosario Toscano. Heuristic Kalman Algorithm{[}M{]}. Germany:
Springer London: 2013, pp: 34-45.

\protect\phantomsection\label{_Toc9629}{}\textbf{致谢}

正文文字正文文字正文文字正文文字正文文字正文文字正文文字正文文字正文文字正文文字正文文字正文文字正文文字正文文字正文文字正文文字正文文字正文文字正文文字正文文字正文文字正文文字正文文字正文文字。

正文文字正文文字正文文字正文文字正文文字正文文字正文文字正文文字正文文字正文文字正文文字正文文字正文文字正文文字正文文字正文文字正文文字正文文字正文文字正文文字正文文字正文文字正文文字正文文字。

\textbf{\hfill\break
}\protect\phantomsection\label{_Toc32605}{}\textbf{附录}

\end{document}
